\chapter{Methodology}
\label{cha:2}

\section{Surface processing pipeline}

\begin{itemize}
	\item Cad model is input
	\item extraction surfaces 
	\item filtering out irrelevant geometry?
	\item computation of geometric descriptors (normals, curvature, area etc)
	\end{itemize}
	
	
	
	\begin{figure}
		\centering
		\includegraphics[scale=0.35]{images/NN visualization.jpg}
		\caption{Feedforward network visualization with two hidden layers}
		\label{fig:NNvisualization}
	\end{figure}
	
	
	
	
\begin{figure}
	\centering
	\includegraphics[scale=0.3]{images/pipeline faces.jpg}
	\caption{Pipeline for surface processing}
	\label{fig:surfacepipeline}
\end{figure}


\section{Discretized Geometry Processing}
\begin{figure}
	\centering
	\includegraphics[scale=0.3]{images/pipeline discretized.jpg}
	\caption{Discretized Pipeline}
	\label{fig:discretized pipeline}
\end{figure}

\begin{itemize}
	\item conversion of brep surfaces into discretized mesh
	\item adaptive mesh refinement
	\item computationally much heavier
	\item physical contraints into loss function: net clamping force zero, patch size constraint to prevent isloated micro triangles
	\item more uncertainty
\end{itemize}
\begin{figure}
	\centering
	\includegraphics[scale=0.7]{images/adaptivemeshrefinement.png}
	\caption{Adaptive mesh refinement}
	\label{fig:adaptive mesh }
\end{figure}

\section{Scaling model}

\begin{itemize}
	\item HPC cluster (reference them as they ask on their website)
	\item optimize model to run on GPU node (parallelization)
	\item parametric modeling to create a large number of parts
	\item used Blender with base part and modifications such as rotation scaling holes chamfers etc
\end{itemize}


\begin{figure}
	\centering
	\includegraphics[scale=0.5]{images/Parametric modelling.png}
	\caption{Parametrically generating parts}
	\label{fig:parametric modelling}
\end{figure}



\section{Feature enrichment}
\begin{itemize}
\item goal is to capture more geometric information to feed into model, therefore the following features were added
\end{itemize}

\subsection{Outerness}

\begin{figure}
	\centering
	\includegraphics[scale=0.7]{images/outerness.png}
	\caption{Outerness}
	\label{fig:outerness}
\end{figure}

\subsection{Opposing surface validation}

\begin{figure}
	\centering
	\includegraphics[scale=0.7]{images/opposing surface validation.jpg}
	\caption{Opposing surface validation}
	\label{fig:opposing surface validation}
\end{figure}


\subsection{Wall thickness}


\begin{figure}
	\centering
	\includegraphics[scale=0.7]{images/wall thickness.png}
	\caption{Wall Thickness}
	\label{fig:wall thickness}
\end{figure}


\section{Post processing}
\begin{itemize}
	\item AI model outputs a single colored region containing multiple solutions
	\item goal is to extract the different solutions and rank them based on feasibility, reachability, stability etcand provide a visualization
	\item The goal is not to automatically choose the best solution, it is a tool to help the process planner or operator by providing a suggestion which he/she should approve
\end{itemize}


\begin{figure}
	\centering
	\includegraphics[scale=0.3]{images/postprocessing.jpg}
	\caption{Visualization of different solutions on the same part in post processing}
	\label{fig:postprocessing}
\end{figure}


\section{Conclusion}

